% Fichier complété par tools/complete_missing_corrections.py.
% Source : CIN/CIN-02-VitesseAcceleration/18_Maxpid/18_Maxpid.tex
% Statut : rédigé (2 question(s)).
\begin{corrigebox}[Corrigé rédigé]
\CorrigeQuestion{1}
La vitesse du point $G$ du solide 4 s'obtient par composition le long
                de la chaîne cinématique. Au point $G$ :
                $\{\mathcal V(4/0)\}_G=\{\vec\Omega_{4/0};\vec V(G,4/0)\}_G$, avec
                $\vec\Omega_{4/0}$ égal à la somme algébrique des vitesses angulaires des
                pivots traversés et $\vec V(G,4/0)$ obtenu par changement de point.
\CorrigeQuestion{2}
On dérive la vitesse dans $\mathcal R_0$ :
                $\vec\Gamma(G,4/0)=d\vec V(G,4/0)/dt|_0$. Chaque barre apporte un terme
                tangentiel $\ddot q\,\vec k_0\wedge\vec r$ et un terme centripète
                $\dot q\,\vec k_0\wedge(\dot q\,\vec k_0\wedge\vec r)$; la glissière
                apporte $\ddot\lambda$ suivant son axe et, si l'axe tourne, les termes de
                Coriolis $2\dot\lambda\vec\Omega\wedge\vec e$.
\end{corrigebox}
